% DEFINED:   sec:conv-unit, sec:weight-sharing, sec:conv-params, eq:conv-params
% RESOLVED:  ch:foundations, sec:mlp-backprop (backward into Part I)
% FORWARD:   sec:conv-forward (Task 4)

\section{One unit, reused everywhere}\label{sec:conv-unit}

The multilayer perceptron of \Cref{ch:foundations} classifies the $8 \times
8$ UCI digits at about 96\% test accuracy. That looks reasonable until you
ask what the model believes about its input. The MLP treats each of the 64
pixels as an independent feature. Pixel $(3, 4)$ and pixel $(3, 5)$, which
are neighbors in the image, are no more related to it than pixel $(3, 4)$
and pixel $(7, 0)$, which are at opposite ends. The MLP has no concept of
adjacency. It learns relationships between pixels from data, exactly as it
would learn relationships between any other tabular features.

A convolutional unit refuses to start from that blank slate. Recall that a
\texttt{Linear} layer is one or more weight vectors $\mathbf{w}$ paired with
biases $b$, computing a logit $z = \mathbf{w} \cdot \mathbf{x} + b$
(\Cref{sec:mlp-backprop}). A convolutional unit is the same dot product,
restricted to a small \textbf{kernel} of size $k \times k$ and applied at
every position of a $2$-D input:
\[
  z_{r, c} = \sum_{i=0}^{k-1} \sum_{j=0}^{k-1}
             W_{i, j}\, x_{r+i,\, c+j} + b.
\]
The kernel $W$ is shared across all output positions $(r, c)$. There is
\emph{one} $W$, not one per position. That single restriction is the whole
definition of a convolution. The small window is the \textbf{locality}
prior: a unit looks at a neighborhood, never the entire input at once.

Everything else is bookkeeping, and there is little of it:
\begin{itemize}
\item \textbf{Borders.} When the kernel runs off the edge of the input, we
  either drop those positions (no padding, the choice this chapter uses) or
  pad the input with zeros to preserve the spatial size.
\item \textbf{Stride.} Sliding by one pixel between applications is the
  default used here; a larger stride downsamples the output.
\item \textbf{Channels.} A real image has several input channels (RGB has
  three), and a layer has several output channels too, one per independent
  kernel run in parallel. With $C_{\text{in}}$ input channels and
  $C_{\text{out}}$ output channels, a kernel has shape $(C_{\text{in}}, k,
  k)$, dotting through every input channel at each spatial position, and the
  layer stacks $C_{\text{out}}$ such kernels.
\end{itemize}

For single-channel $8 \times 8$ digits with $k = 3$, no padding, stride one,
and four output channels, the layer produces an output of shape $(4, 6, 6)$:
four feature maps, each $6 \times 6$ because the kernel cannot fit at the
last two positions in either dimension. The forward pass that computes it is
shown in \Cref{sec:conv-forward}.

\section{Weight sharing is the translation prior}\label{sec:weight-sharing}

Reusing $W$ at every position is what makes the layer \textbf{translation
equivariant}. Stated cleanly: if $x'$ is $x$ shifted by $(\Delta r, \Delta
c)$, then the layer's output on $x'$ is its output on $x$ shifted by the same
$(\Delta r, \Delta c)$, modulo whatever falls off the borders. A detector
that fires on a stroke in one corner fires on the same stroke anywhere.

This is not a property the network \emph{learns}. It is a property of the
\emph{operator}: true at initialization, true after training, true for every
kernel. The reason is the sharing itself. Shift the input by one pixel and
every windowed dot product is now computed at a position also shifted by one,
so the output simply slides. The MLP has no such guarantee. Each of its
output units carries its own weight vector tied to specific coordinates of
the input, so translating the input scrambles which weights see which pixels.

Weight sharing and translation equivariance are therefore the same fact, read
once from the parameter side and once from the function-symmetry side. The
architecture commits to a symmetry of the data by reusing parameters across
the positions that symmetry relates.

\section{The parameter count tells the story}\label{sec:conv-params}

The compactness of weight sharing shows up immediately in the parameter
count, and the count is the cleanest single statement of the prior.

The MLP for the $8 \times 8$ digits has shape $64 \to 32 \to 10$
(\Cref{sec:mlp-backprop}). Its parameter count is
\[
  64 \cdot 32 + 32 + 32 \cdot 10 + 10 = 2410.
\]
A small convolutional network for the same task uses one conv layer of four
$3 \times 3$ kernels followed by a fully connected head:
\begin{lstlisting}
Conv2D(in_channels=1, out_channels=4, kernel_size=3)
ReLU()
Linear(144, 10)
\end{lstlisting}
The conv layer has $4 \cdot (1 \cdot 9) + 4 = 40$ parameters. With no padding
and stride one its output is $4 \times 6 \times 6 = 144$ units, so the head is
\texttt{Linear(144, 10)} with $144 \cdot 10 + 10 = 1450$ parameters. The total
is $40 + 1450 = 1490$, about 38\% smaller than the MLP. Almost all of the
model's spatial capacity sits in the 40 parameters of the conv kernel; the
head is a thin reader on top of the feature maps.

The comparison is sharper if we ask what it would cost to replace the conv
layer with a \texttt{Linear} layer producing the same 144-dimensional output.
That layer is \texttt{Linear(64, 144)}, with $64 \cdot 144 + 144 = 9360$
parameters. The convolution accomplishes the same input-to-output shape
transformation with 40 parameters instead of 9360:
\begin{equation}\label{eq:conv-params}
  \frac{9360}{40} = 234.
\end{equation}
That factor of 234 is the explicit price of \emph{not} committing to locality
and translation equivariance. The forward cost is modest in proportion: four
output channels, one input channel, a $3 \times 3$ kernel, over $6 \times 6$
output positions is $4 \cdot 1 \cdot 9 \cdot 36 = 1296$ multiply-adds per
example. The prior buys a $234\times$ reduction in parameters at the cost of a
fixed, transparent amount of arithmetic. Both numbers are regenerated by the
chapter's paired notebook.
